\section{Screening Protocol}\label{sec:signatures}
\subsection{Mechanism Classes and Probes}
The classes are a taxonomy by the force a perturbation adds to an agent at a rest state. A \emph{Prevent} mechanism adds none: it acts only while the crowd moves, changes how often deadlocks form and leaves a snapshot held. A \emph{Reshape} mechanism adds a deterministic force that moves or removes the rest state (Sect.~\ref{sec:rest}). A \emph{Heat} mechanism adds a random force of zero mean and carries the state over a barrier $\Delta U$ by fluctuations. Activated escape, at the Kramers rate set by $\beta$ (Eq.~\eqref{eq:teff}), is the high-barrier limit of Heat, tested here at $\beta\ge1.5$. The three forces exclude each other, and so do the classes. Reshape has two subtypes: an axial one is goal-directed, and a lateral one is handed and removes a strict minimum once $|b|$ exceeds a threshold $b_f$ (Sect.~\ref{sec:pair}). A composite perturbation is classed by its components, measured on their own.

A \emph{probe} is a quench variant to which the classes respond differently. A replay displaces every coordinate by $10^{-9}$\,m. The time-step probe repeats a quench at $\Delta t=0.005$ and 0.02\,s with per-step or $\sqrt{\Delta t}$-scaled noise (Sect.~\ref{sec:teff}). The \emph{rest-drive push} replaces the ORCA term by its drive part at rest, $m_iv_{0,i}/\tau_{\rm orca}$, directed towards ($+$) or away from ($-$) the door. The bias probe applies $\pm b$, and at a standoff the first exit records which member leaves first. The corrected-implementation probe swaps the legacy for the RVO2 solver.

Each class predicts its own response to each probe. Under a Prevent mechanism a snapshot stays held in every probe, and neither the push nor the bias changes it. An axial Reshape mechanism frees the snapshot with identical replays and a freed fraction that does not depend on $\Delta t$; it frees under the push towards the door and not under the push away, and at a standoff no member is favoured. A lateral Reshape mechanism also frees with identical replays and a flat $\Delta t$ response, but only past $b_f$, with either sign of $b$, and the member on the side favoured by $b$ exits first. A Heat mechanism frees with exponential times to free at $\beta\ge1.5$; its noise seeds disagree, its freed fraction falls as $\Delta t$ grows under per-step noise and stays flat under $\sqrt{\Delta t}$-scaled noise, and no member is favoured. The outcomes of all three classes persist in the corrected implementation. A mechanism of any class carries the \emph{defect flag} when outcomes it produces vanish in the corrected implementation, and \emph{replay sensitivity} when a replay changes its outcome.

Some responses hold by construction: every Prevent response, the disagreement of noise seeds, the per-step $\Delta t$ response of Heat and the favoured side of a lateral bias. The others are empirical: time to free, the final state of a whole run, the overlap of unresolved sets, the first exit under a deterministic resolver and the solver swap.

\subsection{Matching and Separation}\label{sec:protocol}
\begin{figure*}[t]\centering% Fig. (protocol): the five steps of the screening protocol as a greyscale vector flowchart, full \textwidth.
% Input from paper/sections/04_protocol.tex inside a figure* environment; needs tikz with the libraries
% arrows.meta and calc (both loaded in paper/main.tex). No fills: black text and arrows, grey frames.
% One column per step, left to right; terms are defined in the glossary and Sect. 4, so the boxes carry lists,
% not sentences. Classes after ruling R6: Prevent; Reshape, axial or lateral; Heat; the defect flag is a
% separate attribute in its own dashed box. Text is \scriptsize = 7 pt in llncs.
\begin{tikzpicture}[
    font=\scriptsize,
    every node/.style={outer sep=0pt},
    pbox/.style={anchor=north west, inner sep=2.5pt},
    frame/.style={draw=black!60, line width=0.5pt},
    flow/.style={-{Stealth[length=1.4mm,width=1.1mm]}, line width=0.6pt, color=black},
  ]
  \pgfmathsetlengthmacro{\gap}{0.30cm}%
  % column text widths (sum + 5 x 5pt inner sep + 4 gaps = \textwidth)
  \pgfmathsetlengthmacro{\pwa}{0.175*(\textwidth-4*\gap-25.5pt)}%
  \pgfmathsetlengthmacro{\pwb}{0.220*(\textwidth-4*\gap-25.5pt)}%
  \pgfmathsetlengthmacro{\pwc}{0.225*(\textwidth-4*\gap-25.5pt)}%
  \pgfmathsetlengthmacro{\pwd}{0.175*(\textwidth-4*\gap-25.5pt)}%
  \pgfmathsetlengthmacro{\pwe}{0.205*(\textwidth-4*\gap-25.5pt)}%
  % cell text: ragged right, no hyphenation; every box takes the height of the tallest step
  \def\pfmt{\raggedright\hyphenpenalty=10000\exhyphenpenalty=10000 }%
  \def\pA{\textbf{1 Snapshots}\\[1pt]saved rest states:\\ lone keystones\\ second traps\\ pair standoffs}%
  \def\pB{\textbf{2 Probes}\\[1pt]quench\\ replay ($10^{\textstyle -9}$\,m)\\ time step $\Delta t$ with per-step or $\sqrt{\Delta t}$-scaled noise\\ rest-drive push~$\pm$\\ lateral bias $\pm b$\\ RVO2 solver}%
  \def\pC{\textbf{3 Responses}\\[1pt]freed fraction at~$T$\\ times to free\\ $\Delta t$ dependence\\ handedness under bias\\ replay sensitivity\\ (paired by snapshot)}%
  \def\pD{\textbf{4 Class}\\[1pt]Prevent\\ Reshape, axial or lateral\\ Heat}%
  \def\pF{\textbf{Defect flag}\\[1pt]outcome vanishes with the RVO2 solver}%
  \def\pE{\textbf{5 Compare}\\[1pt]resolver against its matched surrogate\\ same snapshots\\ paired tests, Holm-corrected}%
  % text nodes first, frames afterwards down to the lowest step box (the class box keeps its own height)
  \node[pbox] (s1) at (0,0) {\parbox[t]{\pwa}{\pfmt\pA}};
  \node[pbox] (s2) at ($(s1.north east)+(\gap,0)$) {\parbox[t]{\pwb}{\pfmt\pB}};
  \node[pbox] (s3) at ($(s2.north east)+(\gap,0)$) {\parbox[t]{\pwc}{\pfmt\pC}};
  \node[pbox] (s4) at ($(s3.north east)+(\gap,0)$) {\parbox[t]{\pwd}{\pfmt\pD}};
  \node[pbox] (fl) at ($(s4.south west)+(0,-\gap)$) {\parbox[t]{\pwd}{\pfmt\pF}};
  \node[pbox] (s5) at ($(s4.north east)+(\gap,0)$) {\parbox[t]{\pwe}{\pfmt\pE}};
  \global\dimen1=0pt
  \foreach \n in {s1,s2,s3,s5}{\pgfpointanchor{\n}{south}\pgfgetlastxy{\px}{\py}\ifdim\py<\dimen1\global\dimen1=\py\fi}%
  \foreach \n in {s1,s2,s3,s5}{\draw[frame] (\n.north west) rectangle (\n.north east |- 0pt,\the\dimen1);}%
  \draw[frame] (s4.north west) rectangle (s4.south east);%
  \draw[frame, dashed] (fl.north west) rectangle (fl.south east);%
  \foreach \a/\b in {s1/s2, s2/s3, s3/s4, s4/s5}{\draw[flow] ($(\a.north east)+(0,-0.45cm)$) -- ($(\b.north west)+(0,-0.45cm)$);}%
  \coordinate (pbot) at (0pt,\the\dimen1); \coordinate (s3r) at ($(s3.north east)+(-0.35cm,0)$);%
  \draw[flow] (s3r |- pbot) |- (fl.west);%
\end{tikzpicture}%
\caption{Screening protocol. Snapshots are probed; the responses assign a class and the defect flag, and the resolver is compared with its matched surrogate on the same snapshots.}\label{fig:protocol}\end{figure*}
The protocol (Fig.~\ref{fig:protocol}) compares a \emph{resolver}, the perturbation credited with resolution, with a \emph{surrogate} of another class matched to it at a horizon $T$. The two are matched when the unadjusted 95\% interval of the paired difference of their freed fractions at $T$ (completed fractions for whole runs) covers zero; the Holm-adjusted $p$ of the same test is reported separately. A match therefore holds for a given horizon. A stricter test asks for equivalence within $\pm\RvAfiveMargin$ in freed fraction, the 95\% resolution of force-coupled ORCA's own freed fraction on \EtwoNLK--\PairN\ snapshots: two one-sided tests at level 0.05~\cite{schuirmann1987} declare equivalence when the 90\% interval of the difference lies inside this margin.

Quenches run to $T=120$\,s and whole runs to 600\,s. Force-coupled ORCA frees its last screening-set blocker by \RevAeightSTRvoLatest\,s and completes its last whole run by \RevAeightWholeRvoLatest\,s, but noise keeps freeing after these times, so a match with noise holds only from some horizon on. The two criteria give two such horizons, the first from which equality is not rejected and a later one from which equivalence holds; we report the second.

A probe separates two perturbations when a paired test on the same snapshots or spawn seeds rejects equality at level 0.05 after Holm's correction (Sect.~\ref{sec:stats}); the overlap $\Omega$ takes a hypergeometric test.
